\documentclass[10pt,a4paper]{amsart}
\usepackage{amsmath,amssymb,amsthm,graphicx}
% Retained minimal preamble for the Stacks Project chapter more-morphisms.tex (More on Morphisms).
% Verbatim excerpt of the chapter preamble: Xy-pic setup (native preamble lines 14-18) and the native
% theorem declarations with the mathematical macros (native preamble lines 154-217). Copied mechanically;
% nothing added, removed or redefined. The wrapper class and the amsmath/amssymb/amsthm/graphicx base
% packages are supplied by the assembler, so the class/externaldocument part of the native preamble is omitted.
% For commutative diagrams we use Xy-pic
\usepackage[all]{xy}

% We use 2cell for 2-commutative diagrams.
\xyoption{2cell}
% Theorem environments.
%
\theoremstyle{plain}
\newtheorem{theorem}[subsection]{Theorem}
\newtheorem{proposition}[subsection]{Proposition}
\newtheorem{lemma}[subsection]{Lemma}

\theoremstyle{definition}
\newtheorem{definition}[subsection]{Definition}
\newtheorem{example}[subsection]{Example}
\newtheorem{exercise}[subsection]{Exercise}
\newtheorem{situation}[subsection]{Situation}

\theoremstyle{remark}
\newtheorem{remark}[subsection]{Remark}
\newtheorem{remarks}[subsection]{Remarks}

\numberwithin{equation}{subsection}

% Macros
%
\def\lim{\mathop{\mathrm{lim}}\nolimits}
\def\colim{\mathop{\mathrm{colim}}\nolimits}
\def\Spec{\mathop{\mathrm{Spec}}}
\def\Hom{\mathop{\mathrm{Hom}}\nolimits}
\def\Ext{\mathop{\mathrm{Ext}}\nolimits}
\def\SheafHom{\mathop{\mathcal{H}\!\mathit{om}}\nolimits}
\def\SheafExt{\mathop{\mathcal{E}\!\mathit{xt}}\nolimits}
\def\Sch{\mathit{Sch}}
\def\Mor{\mathop{\mathrm{Mor}}\nolimits}
\def\Ob{\mathop{\mathrm{Ob}}\nolimits}
\def\Sh{\mathop{\mathit{Sh}}\nolimits}
\def\NL{\mathop{N\!L}\nolimits}
\def\CH{\mathop{\mathrm{CH}}\nolimits}
\def\proetale{{pro\text{-}\acute{e}tale}}
\def\etale{{\acute{e}tale}}
\def\QCoh{\mathit{QCoh}}
\def\Ker{\mathop{\mathrm{Ker}}}
\def\Im{\mathop{\mathrm{Im}}}
\def\Coker{\mathop{\mathrm{Coker}}}
\def\Coim{\mathop{\mathrm{Coim}}}

% Boxtimes
%
\DeclareMathSymbol{\boxtimes}{\mathbin}{AMSa}{"02}

%
% Macros for moduli stacks/spaces
%
\def\QCohstack{\mathcal{QC}\!\mathit{oh}}
\def\Cohstack{\mathcal{C}\!\mathit{oh}}
\def\Spacesstack{\mathcal{S}\!\mathit{paces}}
\def\Quotfunctor{\mathrm{Quot}}
\def\Hilbfunctor{\mathrm{Hilb}}
\def\Curvesstack{\mathcal{C}\!\mathit{urves}}
\def\Polarizedstack{\mathcal{P}\!\mathit{olarized}}
\def\Complexesstack{\mathcal{C}\!\mathit{omplexes}}
% \Pic is the operator that assigns to X its picard group, usage \Pic(X)
% \Picardstack_{X/B} denotes the Picard stack of X over B
% \Picardfunctor_{X/B} denotes the Picard functor of X over B
\def\Pic{\mathop{\mathrm{Pic}}\nolimits}
\def\Picardstack{\mathcal{P}\!\mathit{ic}}
\def\Picardfunctor{\mathrm{Pic}}
\def\Deformationcategory{\mathcal{D}\!\mathit{ef}}
\usepackage{cleveref}
\crefname{originalreference}{original reference}{original references}
\Crefname{originalreference}{Original reference}{Original references}
\makeatletter
\newlabel{algebra-lemma-finite-projective}{{00NX}{source}{}{sourcepointer-1}{}}
\newlabel{algebra-lemma-finite-projective@cref}{{[originalreference][0][]00NX}{[1][0][]source}{}{}{}}
\makeatother
\crefname{originalreference}{original reference}{original references}
\Crefname{originalreference}{Original reference}{Original references}
\makeatletter
\newlabel{algebra-lemma-finite-projective-descends}{{00O1}{source}{}{sourcepointer-2}{}}
\newlabel{algebra-lemma-finite-projective-descends@cref}{{[originalreference][0][]00O1}{[1][0][]source}{}{}{}}
\makeatother
\crefname{originalreference}{original reference}{original references}
\Crefname{originalreference}{Original reference}{Original references}
\makeatletter
\newlabel{algebra-section-completion-noetherian}{{0BNH}{source}{}{sourcepointer-3}{}}
\newlabel{algebra-section-completion-noetherian@cref}{{[originalreference][0][]0BNH}{[1][0][]source}{}{}{}}
\makeatother
\crefname{originalreference}{original reference}{original references}
\Crefname{originalreference}{Original reference}{Original references}
\makeatletter
\newlabel{coherent-lemma-formal-functions-stalk}{{02OD}{source}{}{sourcepointer-4}{}}
\newlabel{coherent-lemma-formal-functions-stalk@cref}{{[originalreference][0][]02OD}{[1][0][]source}{}{}{}}
\makeatother
\crefname{originalreference}{original reference}{original references}
\Crefname{originalreference}{Original reference}{Original references}
\makeatletter
\newlabel{coherent-proposition-proper-pushforward-coherent}{{02O5}{source}{}{sourcepointer-5}{}}
\newlabel{coherent-proposition-proper-pushforward-coherent@cref}{{[originalreference][0][]02O5}{[1][0][]source}{}{}{}}
\makeatother
\crefname{originalreference}{original reference}{original references}
\Crefname{originalreference}{Original reference}{Original references}
\makeatletter
\newlabel{cohomology-lemma-projection-formula}{{01E8}{source}{}{sourcepointer-6}{}}
\newlabel{cohomology-lemma-projection-formula@cref}{{[originalreference][0][]01E8}{[1][0][]source}{}{}{}}
\makeatother
\crefname{originalreference}{original reference}{original references}
\Crefname{originalreference}{Original reference}{Original references}
\makeatletter
\newlabel{cohomology-proposition-vanishing-Noetherian}{{02UZ}{source}{}{sourcepointer-7}{}}
\newlabel{cohomology-proposition-vanishing-Noetherian@cref}{{[originalreference][0][]02UZ}{[1][0][]source}{}{}{}}
\makeatother
\crefname{originallemma}{original lemma}{original lemmas}
\Crefname{originallemma}{Original lemma}{Original lemmas}
\makeatletter
\newlabel{lemma-picard-group-first-order-thickening}{{0C6R}{source}{}{sourcepointer-8}{}}
\newlabel{lemma-picard-group-first-order-thickening@cref}{{[originallemma][0][]0C6R}{[1][0][]source}{}{}{}}
\makeatother
\title{Pullback on Picard groups is a bijection onto the subgroup of fibre-trivial line bundles}
\author{The Stacks Project Authors}
\date{stacks-project:more-morphisms.tex:local-snapshot-20261011; GNU FDL 1.2 or later}
\begin{document}
\maketitle
\noindent\textit{Editorial note.} Counted span is the author's own proof environment, source lines 21357-21411, immediately following the statement at 21345-21355.\par\medskip
\noindent\textit{Editorial note.} Seven cross-chapter pointers (Algebra, Cohomology of Schemes, Coherent, Cohomology) plus the chapter's own first-order-thickening lemma are recorded in assembly.external\_references.\par\medskip
\noindent\textit{Editorial note.} The statement's condition on fibres is used exactly through the vanishing hypothesis; no material is imported from outside the counted proof except as pointers.\par\medskip
\noindent\textit{Allowed standard prerequisites (editorial).}\begin{itemize}
\item ThePicard group of a first order thickening. chapter lemma lemma-picard-group-first-order-thickening, cited in the proof. Provides the infinitesimal step that drives the induction on thickenings.
\item Finiteness and descent of finite projective modules. Algebra, lemmas algebra-lemma-finite-projective and algebra-lemma-finite-projective-descends, cited in the proof. Keeps the relevant modules finite projective and lets them be descended along the thickening. 
\item Cohomology and base change for a proper morphism. Cohomology of Schemes, proposition cohomology-proposition-vanishing-Noetherian and lemma cohomology-lemma-projection-formula, cited in the proof. Vanishing of the first cohomology in the setup and the projection formula for the pushforward of line bundles.
\item Formal functions on the stalk and the completion section. Cohomology of Schemes, lemma coherent-lemma-formal-functions-stalk with Algebra, section algebra-section-completion-noetherian, cited in the proof. Controls the completion of the structure sheaf used in the deformation argument.
\end{itemize}\medskip
\ifdefined\hypertarget\hypertarget{sourcepointer-1}{}\fi
\noindent\textit{Source reference.} Original reference 00NX: focused-seeds/geometry/S9-morph-bijection-on-pic/ref-tags.txt, source lines 4--4. Navigation pointer only: the counted proof cites this statement by \textbackslash{}ref as a standing input; the deduction the counted proof has to make is carried out inside the counted span, and no part of that deduction is delegated to this pointer.\par\medskip
\ifdefined\hypertarget\hypertarget{sourcepointer-2}{}\fi
\noindent\textit{Source reference.} Original reference 00O1: focused-seeds/geometry/S9-morph-bijection-on-pic/ref-tags.txt, source lines 5--5. Navigation pointer only: the counted proof cites this statement by \textbackslash{}ref as a standing input; the deduction the counted proof has to make is carried out inside the counted span, and no part of that deduction is delegated to this pointer.\par\medskip
\ifdefined\hypertarget\hypertarget{sourcepointer-3}{}\fi
\noindent\textit{Source reference.} Original reference 0BNH: focused-seeds/geometry/S9-morph-bijection-on-pic/ref-tags.txt, source lines 6--6. Navigation pointer only: the counted proof cites this statement by \textbackslash{}ref as a standing input; the deduction the counted proof has to make is carried out inside the counted span, and no part of that deduction is delegated to this pointer.\par\medskip
\ifdefined\hypertarget\hypertarget{sourcepointer-4}{}\fi
\noindent\textit{Source reference.} Original reference 02OD: focused-seeds/geometry/S9-morph-bijection-on-pic/ref-tags.txt, source lines 7--7. Navigation pointer only: the counted proof cites this statement by \textbackslash{}ref as a standing input; the deduction the counted proof has to make is carried out inside the counted span, and no part of that deduction is delegated to this pointer.\par\medskip
\ifdefined\hypertarget\hypertarget{sourcepointer-5}{}\fi
\noindent\textit{Source reference.} Original reference 02O5: focused-seeds/geometry/S9-morph-bijection-on-pic/ref-tags.txt, source lines 8--8. Navigation pointer only: the counted proof cites this statement by \textbackslash{}ref as a standing input; the deduction the counted proof has to make is carried out inside the counted span, and no part of that deduction is delegated to this pointer.\par\medskip
\ifdefined\hypertarget\hypertarget{sourcepointer-6}{}\fi
\noindent\textit{Source reference.} Original reference 01E8: focused-seeds/geometry/S9-morph-bijection-on-pic/ref-tags.txt, source lines 9--9. Navigation pointer only: the counted proof cites this statement by \textbackslash{}ref as a standing input; the deduction the counted proof has to make is carried out inside the counted span, and no part of that deduction is delegated to this pointer.\par\medskip
\ifdefined\hypertarget\hypertarget{sourcepointer-7}{}\fi
\noindent\textit{Source reference.} Original reference 02UZ: focused-seeds/geometry/S9-morph-bijection-on-pic/ref-tags.txt, source lines 10--10. Navigation pointer only: the counted proof cites this statement by \textbackslash{}ref as a standing input; the deduction the counted proof has to make is carried out inside the counted span, and no part of that deduction is delegated to this pointer.\par\medskip
\ifdefined\hypertarget\hypertarget{sourcepointer-8}{}\fi
\noindent\textit{Source reference.} Original lemma 0C6R: sources/stacks-more-morphisms/pinned-original-source.tex, source lines 599--599. Navigation pointer only: the counted proof cites this statement by \textbackslash{}ref as a standing input; the deduction the counted proof has to make is carried out inside the counted span, and no part of that deduction is delegated to this pointer.\par\medskip
\begin{lemma}
\label{lemma-bijection-on-Pic}
Let $f : X \to Y$ be a proper morphism of Noetherian schemes
such that $f_*\mathcal{O}_X = \mathcal{O}_Y$, such that
the fibres of $f$ have dimension $\leq 1$, and such that
$H^1(X_y, \mathcal{O}_{X_y}) = 0$ for $y \in Y$.
Then $f^* : \Pic(Y) \to \Pic(X)$ is a bijection onto
the subgroup of $\mathcal{L} \in \Pic(X)$ with
$\mathcal{L}|_{X_y} \cong \mathcal{O}_{X_y}$
for all $y \in Y$.
\end{lemma}

\begin{proof}
By the projection formula
(Cohomology, Lemma \ref{cohomology-lemma-projection-formula})
we see that $f_*f^*\mathcal{N} \cong \mathcal{N}$ for
$\mathcal{N} \in \Pic(Y)$.
We claim that for $\mathcal{L} \in \Pic(X)$ with
$\mathcal{L}|_{X_y} \cong \mathcal{O}_{X_y}$
for all $y \in Y$ we have $\mathcal{N} = f_*\mathcal{L}$
is invertible and $\mathcal{L} \cong f^*\mathcal{N}$.
This will finish the proof.

\medskip\noindent
The $\mathcal{O}_Y$-module $\mathcal{N} = f_*\mathcal{L}$
is coherent by Cohomology of Schemes, Proposition
\ref{coherent-proposition-proper-pushforward-coherent}.
Thus to see that it is an invertible $\mathcal{O}_Y$-module,
it suffices to check on stalks (Algebra, Lemma
\ref{algebra-lemma-finite-projective}).
Since the map from a Noetherian local ring to its completion
is faithfully flat, it suffices to check the completion
$(f_*\mathcal{L})_y^\wedge$ is free (see
Algebra, Section \ref{algebra-section-completion-noetherian} and
Lemma \ref{algebra-lemma-finite-projective-descends}).
For this we will use the theorem of formal functions as formulated in
Cohomology of Schemes, Lemma \ref{coherent-lemma-formal-functions-stalk}.
Since $f_*\mathcal{O}_X = \mathcal{O}_Y$ and hence
$(f_*\mathcal{O}_X)_y^\wedge \cong \mathcal{O}_{Y, y}^\wedge$,
it suffices to show that $\mathcal{L}|_{X_n} \cong \mathcal{O}_{X_n}$
for each $n$ (compatibly for varying $n$.
By Lemma \ref{lemma-picard-group-first-order-thickening}
we have an exact sequence
$$
H^1(X_y, \mathfrak m_y^n\mathcal{O}_X/\mathfrak m_y^{n + 1}\mathcal{O}_X)
\to \Pic(X_{n + 1}) \to \Pic(X_n)
$$
with notation as in the theorem on formal functions.
Observe that we have a surjection
$$
\mathcal{O}_{X_y}^{\oplus r_n} \cong
\mathfrak m_y^n/\mathfrak m_y^{n + 1} \otimes_{\kappa(y)} \mathcal{O}_{X_y}
\longrightarrow
\mathfrak m_y^n\mathcal{O}_X/\mathfrak m_y^{n + 1}\mathcal{O}_X
$$
for some integers $r_n \geq 0$.
Since $\dim(X_y) \leq 1$ this surjection induces a surjection
on first cohomology groups (by the vanishing of cohomology in degrees $\geq 2$
coming from Cohomology, Proposition
\ref{cohomology-proposition-vanishing-Noetherian}).
Hence the $H^1$ in the sequence is zero and the transition maps
$\Pic(X_{n + 1}) \to \Pic(X_n)$ are injective as desired.

\medskip\noindent
We still have to show that $f^*\mathcal{N} \cong \mathcal{L}$.
This is proved by the same method and we omit the details.
\end{proof}

\end{document}
